\hmsubsection{Hardware Integration}{UMA}{}\label{sec:hardware-integration}
This section describes how the hardware components introduced in 
Section~\ref{sec:hardware-architecture} are physically integrated into a working system. The 
description covers the data and power connections between the host 
processor, the motor controller, the camera, and the actuators, the 
configuration of the Dynamixel servos, and the mechanical and runtime 
measures considered to reduce thermal loading during prolonged 
operation.

\hmsubsubsection{Physical Connections}{UMA}{}
The Raspberry Pi~5 acts as the central host of the system and is 
connected to two peripheral devices through USB. The OAK-D~S2 
camera connects directly to one of the Pi's USB~3.0 ports through 
a USB-A to USB-C cable, providing both data transfer and the 
camera's required 5~V power. The OpenRB-150 motor controller connects to a second USB~3.0 port 
through a USB-C to USB-C cable, which carries serial data at 
115\,200~bps. The 12~V power supply for the 
motors is delivered to the OpenRB-150 through its dedicated motor 
power input, separately from the USB connection to the host. The 
Raspberry Pi itself is powered through a third USB-C cable from a 
regulated 5~V supply. All three power and data paths share a 
common ground reference.

The five Dynamixel servos are connected to the OpenRB-150 through 
the two-branch star topology described in Section~\ref{sec:power-architecture} 
(Figure~\ref{fig:power_topology}). Each branch leaves the controller 
through a separate 3-pin TTL cable, and the servos within each 
branch are linked through short JST cables in a daisy chain. The 
TTL bus wiring uses the same physical connector for power and data. 
Each 3-pin Dynamixel TTL connection carries motor supply voltage, 
ground, and a single-ended half-duplex data line on the same connector, 
in accordance with the Dynamixel Protocol 2.0 specification 
\cite{robotis_dynamixel}.

\hmsubsubsection{Motor Configuration}{UMA}{}\label{sec:motor-config}
Before the servos can operate as a coordinated arm, each motor 
must be configured individually with a unique identifier on the 
TTL bus, a fixed baud rate, and a defined zero position. This 
configuration was performed using the ROBOTIS Dynamixel Wizard 
2.0 software with each motor connected to the OpenRB-150 in 
isolation. The five motors were assigned identifiers 1 through 5, 
corresponding to the base, shoulder, elbow, wrist, and gripper 
joints respectively. All five motors were set to a baud rate of 115\,200~bps, matching 
the baud rate used on the USB serial connection between the host 
processor and the OpenRB-150. Using a single rate on both sides 
of the controller simplifies the firmware and ensures that the 
limiting factor in the command-response cycle is the motor's 
mechanical response time rather than any rate mismatch between 
the two buses. While the Dynamixel servos support baud rates up 
to 4.5~Mbps, 115\,200~bps provides sufficient throughput for the 
movement command rate produced by the host, since the system 
issues at most a few commands per second during a typical 
pick-and-place cycle. Each motor 
was rotated to its mechanical centre and the corresponding encoder 
value was registered as the zero position, after which the joint 
limits described in Section~\ref{sec:joint-limits} were applied to prevent 
the motors from rotating into mechanical hard stops during 
operation.

Two of the motors used in the project initially exhibited 
hardware error flags that prevented them from accepting commands. 
The flags were investigated by reading the hardware error status 
register from each motor over the TTL bus and were found to be 
latched overload errors from previous use. A factory reset was 
performed on the affected motors through the Dynamixel Wizard, 
after which the error flags were cleared and the motors operated 
normally. Both motors were retained for use in the final assembly.

\hmsubsubsection{Thermal Protection at Scan Pose}{UMA}{}\label{sec:thermal}
A practical issue identified during integration testing concerned 
the elbow motor (M3, XM430-W210-T). The system spends a significant 
proportion of its operating time holding the arm at a fixed 
\texttt{SCAN\_POSE} configuration, in which the camera is positioned 
above the workspace for ball detection (see Section~\ref{sec:ik-solver}). 
In the initial scan pose, the elbow motor had to continuously resist 
the gravitational torque produced by the forearm, the wrist, the 
gripper, and the camera. Holding torque draws current continuously, 
and the XM430-W210-T was observed to draw approximately 0.47~A while 
stationary at this pose. Over extended periods, this current caused 
the motor to heat up, which could trigger an overheating error flag 
and cause the motor to shut down mid-operation.

To address this, a reduced current limit of 300~mA was
implemented and tested on the elbow motor whenever the arm
was parked at \texttt{SCAN\_POSE}. The intention was to cap
the steady-state current and thereby reduce heat generation
while still providing sufficient torque to hold the arm in
position under nominal load.

\textbf{Outcome.}
Testing showed that the reduced current limit did not
produce a measurable improvement in motor temperature
during extended operation. The thermal load on the elbow
motor is fundamentally determined by the gravitational
torque from the weight of the forearm, camera, and gripper
assembly, which the motor must resist regardless of the
current limit setting. A current limit can only prevent
the motor from drawing excessive current during transient
loads; the steady-state holding current is already below
the stall current and cannot be reduced further without
allowing the arm to droop.

The first mechanical mitigation was to reduce the mass of the
components carried by the outer arm. This improved the situation
somewhat by lowering the gravitational load, but it did not remove
the main cause of heating while the arm remained stationary. The
main mitigation for idle scan-pose heating was instead a revised scan
pose. In this configuration, the L\textsubscript{2} structure rests
mechanically on the L\textsubscript{1} piece while the camera remains
oriented toward the workspace. The contact between the two link
structures carries most of the static load that previously had to be
held by M3, so the elbow motor uses almost no holding torque while the
arm is idle in the scan position.

The revised pose preserved the functional purpose of the scan pose: the
camera still had a clear and sufficiently wide view of the balls in the
workspace. After this change, no problematic heating was observed while
holding the revised scan pose during the available integration tests.
This observation applies specifically to stationary idle time at
\texttt{SCAN\_POSE}; it does not mean that the thermal issue was fully
resolved during complete sorting cycles.

\textbf{Conclusion.}
The thermal protection logic was therefore removed from the
deployed system. The effective mitigation in the deployed design was
the mechanically supported scan pose, supplemented by the earlier mass
reduction work, for reducing idle scan-pose heating. During operational
testing, the robot was able to sort 10 balls consecutively, but the
motors still became somewhat hot during movement. Extended continuous
operation should therefore still be treated as a limitation and
monitored, because the scan-pose change mitigates stationary heating
rather than eliminating movement-related thermal risk.
